% ソフトウェア工学A チーム開発結果
%
% 図の貼り付けと表だけを埋めてある。本文（説明の文章）は
% 「%% TODO」のコメントを置いた箇所に書き足すこと。
%
% コンパイル: xelatex report.tex
% 画像: docs/uml/images/*.png を参照する。Overleaf等へ持っていく場合は
%       images/ を report.tex と同じ階層に置けば graphicspath が拾う。

\documentclass[xelatex,ja=standard,a4paper,12pt]{bxjsarticle}

\usepackage{zxjatype}
\usepackage[haranoaji]{zxjafont}
\usepackage{tabularx}
\usepackage{array}
\usepackage{seqsplit}
\usepackage{float}
\usepackage{graphicx}
\usepackage{multirow}

\graphicspath{{../uml/images/}{images/}{./}}

% 図の貼り付け。ファイル名にアンダースコアが含まれるため \detokenize で囲む。
% #1: ファイル名, #2: \linewidth に対する倍率, #3: キャプション, #4: ラベル
\newcommand{\umlfigure}[4]{%
    \begin{figure}[H]
        \centering
        \includegraphics[width=#2\linewidth]{\detokenize{#1}}
        \caption{#3}
        \label{fig:#4}
    \end{figure}%
}

\begin{document}

\begin{center}
\section*{ソフトウェア工学A チーム開発結果}
\end{center}

\begin{flushright}
    student\_number name
\end{flushright}

\subsection*{1. チーム番号およびメンバ構成}
チーム番号: 28

%% TODO: メンバ構成を書く
メンバ構成

\subsection*{2. チーム開発結果}

\subsubsection*{2.1 ドメイン分析}

%% TODO: 何を分析し、どう表現したかの説明を書く

\umlfigure{01_domain_class_diagram.png}{0.34}{ドメイン分析のPUMLクラス図}{domain-class}

\umlfigure{02_domain_object_diagram.png}{1.0}{ドメイン分析のPUMLオブジェクト図}{domain-object}

%% TODO: オブジェクト図でクラス図を検証した結果を書く
%% （2つの具体的な状況を置き、予約—チェックインの多重度 0..1 の下限0を検証した、など）

\begin{table}[H]
    \centering
    \caption{ドメイン分析のレビュー項目}
    \label{tab:domain-analysis-review}

    \renewcommand{\arraystretch}{1.4}

    \begin{tabularx}{\linewidth}{
        |>{\centering\arraybackslash}p{2.2cm}
        |X
        |>{\centering\arraybackslash}p{1.2cm}|
    }
        \hline
        対象 & \multicolumn{1}{c|}{項目} & 達成度 \\
        \hline

        \multirow{4}{*}{全体}
        & 主要な概念および概念間の関係が網羅されている
        & $\bigcirc$ \\
        \cline{2-3}

        & 対象とする問題領域の用語表・辞書として機能する
        & $\bigcirc$ \\
        \cline{2-3}

        & 具体的な実装方法の推定とはなっていない
        & $\bigcirc$ \\
        \cline{2-3}

        & 配置や色、フォントなどが読みやすい
        & $\bigcirc$ \\
        \hline

        \multirow{2}{*}{クラス}
        & 主に名詞・名詞句・代名詞が、クラス・属性として挙げられている
        & $\bigcirc$ \\
        \cline{2-3}

        & 直観的で明確なクラス名がつけられている
        & $\bigcirc$ \\
        \hline

        \multirow{3}{*}{関係}
        & 全ての関係について、両端を用いて違和感なく読める
        & $\bigcirc$ \\
        \cline{2-3}

        & 汎化・集約の関係について、必然性がある
        & $\bigcirc$ \\
        \cline{2-3}

        & 必要に応じて多重度・ロールが用いられている
        & $\bigcirc$ \\
        \hline
    \end{tabularx}
\end{table}

\subsubsection*{2.2 要求分析}

作成したPUMLユースケース図を図\ref{fig:usecase}に示す。ユースケース図の各ユースケース(部屋を予約する・予約をキャンセルする・チェックインする・チェックアウトする・予約内容を確認する)についてのユースケース記述をそれぞれ説明する。

\umlfigure{03_usecase_diagram.png}{0.75}{要求分析のPUMLユースケース図}{usecase}

\begin{table}[H]
    \centering
    \caption{ユースケース記述: UC1 部屋を予約する}
    \label{tab:uc1}
    \renewcommand{\arraystretch}{1.3}
    \begin{tabularx}{\linewidth}{|>{\centering\arraybackslash}p{2.2cm}|X|}
        \hline
        \multicolumn{1}{|c|}{項目} & \multicolumn{1}{c|}{内容} \\
        \hline
        アクター & 利用者 \\
        \hline
        概要 & 利用者が希望する客室を選び予約を確定し、予約番号を受け取る \\
        \hline
        事前条件 & 利用者が予約対象のホテル・宿泊日・客室タイプを把握している \\
        \hline
        事後条件 & 予約が「予約済み」状態で登録され、予約番号が発行されている \\
        \hline
        基本フロー &
        1. 利用者が宿泊日・客室タイプを指定する \newline
        2. システムが該当する空室一覧を提示する \newline
        3. 利用者が客室を選択し、氏名・連絡先を入力する \newline
        4. システムが予約を登録する \newline
        5. システムが予約番号を発行する \newline
        6. システムが予約番号を利用者に提示する \\
        \hline
        代替フロー &
        3a. 利用者が候補の客室すべてに対し気に入らず再検索を要求する場合 \newline
        \hspace*{1em}3a1. 利用者が宿泊日・客室タイプを変更する \newline
        \hspace*{1em}3a2. 基本フローの手順2に戻る \\
        \hline
        例外フロー &
        2a. システムが該当する空室を1件も検出できない場合 \newline
        \hspace*{1em}2a1. システムが「空室なし」を利用者に通知する \newline
        \hspace*{1em}2a2. ユースケースを終了する \newline
        4a. システムが予約を登録する際、選択した客室が直前に他の予約で確保されていた場合（予約の競合） \newline
        \hspace*{1em}4a1. システムが「該当客室は確保できない」ことを利用者に通知する \newline
        \hspace*{1em}4a2. 基本フローの手順2に戻る \\
        \hline
    \end{tabularx}
\end{table}

\begin{table}[H]
    \centering
    \caption{ユースケース記述: UC2 チェックインする}
    \label{tab:uc2}
    \renewcommand{\arraystretch}{1.3}
    \begin{tabularx}{\linewidth}{|>{\centering\arraybackslash}p{2.2cm}|X|}
        \hline
        \multicolumn{1}{|c|}{項目} & \multicolumn{1}{c|}{内容} \\
        \hline
        アクター & フロント係（利用者はシステムを操作せず、窓口で対応を受ける） \\
        \hline
        概要 & 利用者が予約番号を提示し、フロント係がチェックイン処理を行い部屋番号を通知する \\
        \hline
        事前条件 & 予約が「予約済み」状態で存在する \\
        \hline
        事後条件 & 予約が「チェックイン済み」状態になり、部屋番号が利用者に通知されている \\
        \hline
        基本フロー &
        1. 利用者が予約番号を提示する \newline
        2. フロント係が予約内容を確認する（include UC4） \newline
        3. システムが該当予約の内容をフロント係に提示する \newline
        4. フロント係がチェックイン処理を実行する \newline
        5. システムが予約状態をチェックイン済みに更新する \newline
        6. システムが部屋番号をフロント係に提示する \newline
        7. フロント係が部屋番号を利用者に通知する \\
        \hline
        代替フロー & なし（予約番号の提示を前提とする） \\
        \hline
        例外フロー &
        3a. システムが該当する予約を検出できない場合 \newline
        \hspace*{1em}3a1. システムが「予約なし」をフロント係に通知する \newline
        \hspace*{1em}3a2. フロント係が利用者に再確認を依頼する \newline
        \hspace*{1em}3a3. ユースケースを終了する \\
        \hline
    \end{tabularx}
\end{table}

\begin{table}[H]
    \centering
    \caption{ユースケース記述: UC3 チェックアウトする}
    \label{tab:uc3}
    \renewcommand{\arraystretch}{1.3}
    \begin{tabularx}{\linewidth}{|>{\centering\arraybackslash}p{2.2cm}|X|}
        \hline
        \multicolumn{1}{|c|}{項目} & \multicolumn{1}{c|}{内容} \\
        \hline
        アクター & フロント係（利用者はシステムを操作せず、窓口で対応を受ける） \\
        \hline
        概要 & 宿泊終了後、フロント係がチェックアウト処理を行い宿泊料金を確定し、利用者が支払う \\
        \hline
        事前条件 & 予約が「チェックイン済み」状態である \\
        \hline
        事後条件 & 予約が「チェックアウト済み」状態になり、宿泊料金が支払い済みとして記録される \\
        \hline
        基本フロー &
        1. 利用者がフロントでチェックアウトを申告する \newline
        2. システムが宿泊料金を計算する \newline
        3. システムが宿泊料金をフロント係に提示する \newline
        4. フロント係が利用者に料金を案内する \newline
        5. 利用者が料金を支払う \newline
        6. システムが支払いを記録する \newline
        7. システムが予約状態をチェックアウト済みに更新する \\
        \hline
        代替フロー &
        2a. 追加サービスの利用履歴がある場合 \newline
        \hspace*{1em}2a1. システムが追加料金を宿泊料金に加算する \newline
        \hspace*{1em}2a2. 基本フローの手順3に進む \\
        \hline
        例外フロー &
        5a. 利用者が提示された料金で支払いを拒否する場合 \newline
        \hspace*{1em}5a1. フロント係が利用者と料金内容を再確認する \newline
        \hspace*{1em}5a2. 解決しない場合、ユースケースを終了する（チェックアウト未完了） \\
        \hline
    \end{tabularx}
\end{table}

\begin{table}[H]
    \centering
    \caption{ユースケース記述: UC4 予約内容を確認する}
    \label{tab:uc4}
    \renewcommand{\arraystretch}{1.3}
    \begin{tabularx}{\linewidth}{|>{\centering\arraybackslash}p{2.2cm}|X|}
        \hline
        \multicolumn{1}{|c|}{項目} & \multicolumn{1}{c|}{内容} \\
        \hline
        アクター & 利用者、フロント係 \\
        \hline
        概要 & 予約番号をもとに該当予約の内容（宿泊日・客室タイプ・予約状態）を参照する。UC2・UC5からincludeされる共通機能 \\
        \hline
        事前条件 & 確認対象の予約が登録されている \\
        \hline
        事後条件 & 予約内容が提示されている（予約の状態は変化しない） \\
        \hline
        基本フロー &
        1. アクターが予約番号を入力する \newline
        2. システムが該当予約を検索する \newline
        3. システムが予約内容（宿泊日・客室タイプ・予約状態）を提示する \\
        \hline
        代替フロー & なし（予約番号の入力を前提とする） \\
        \hline
        例外フロー &
        2a. システムが該当する予約を検出できない場合 \newline
        \hspace*{1em}2a1. システムが「予約なし」をアクターに通知する \newline
        \hspace*{1em}2a2. ユースケースを終了する \\
        \hline
    \end{tabularx}
\end{table}

\begin{table}[H]
    \centering
    \caption{ユースケース記述: UC5 予約をキャンセルする}
    \label{tab:uc5}
    \renewcommand{\arraystretch}{1.3}
    \begin{tabularx}{\linewidth}{|>{\centering\arraybackslash}p{2.2cm}|X|}
        \hline
        \multicolumn{1}{|c|}{項目} & \multicolumn{1}{c|}{内容} \\
        \hline
        アクター & 利用者 \\
        \hline
        概要 & 利用者がチェックイン前の予約を取り消し、対象客室を空室に戻す \\
        \hline
        事前条件 & 予約が「予約済み」状態で存在する \\
        \hline
        事後条件 & 予約が「キャンセル済み」状態になり、対象客室が空室に戻っている \\
        \hline
        基本フロー &
        1. 利用者が予約番号を提示する \newline
        2. 利用者が予約内容を確認する（include UC4） \newline
        3. システムが該当予約の内容を提示する \newline
        4. 利用者がキャンセルを指示する \newline
        5. システムがキャンセルの確認を求める \newline
        6. 利用者がキャンセルを確定する \newline
        7. システムが予約状態をキャンセル済みに更新する \newline
        8. システムが対象客室を空室に戻す \newline
        9. システムがキャンセル完了を利用者に通知する \\
        \hline
        代替フロー &
        6a. 利用者がキャンセルを取りやめる場合 \newline
        \hspace*{1em}6a1. システムが予約を「予約済み」のまま維持する \newline
        \hspace*{1em}6a2. ユースケースを終了する \\
        \hline
        例外フロー &
        3a. システムが該当する予約を検出できない場合 \newline
        \hspace*{1em}3a1. システムが「予約なし」を利用者に通知する \newline
        \hspace*{1em}3a2. ユースケースを終了する \newline
        4a. 予約が「予約済み」状態でない場合（チェックイン済み・キャンセル済み等） \newline
        \hspace*{1em}4a1. システムが「キャンセル不可」と理由を利用者に通知する \newline
        \hspace*{1em}4a2. ユースケースを終了する \\
        \hline
    \end{tabularx}
\end{table}

%% TODO: アクティビティ図についての説明を書く（UC1・UC2・UC3の3つを作成した理由など）

\umlfigure{05_activity_diagram_reserve.png}{0.62}{要求分析のPUMLアクティビティ図（UC1 部屋を予約する）}{activity-reserve}

\umlfigure{06_activity_diagram_checkin.png}{0.45}{要求分析のPUMLアクティビティ図（UC2 チェックインする）}{activity-checkin}

\umlfigure{07_activity_diagram_checkout.png}{0.45}{要求分析のPUMLアクティビティ図（UC3 チェックアウトする）}{activity-checkout}

\begin{table}[H]
    \centering
    \caption{要求分析のレビュー項目}
    \label{tab:requirement-analysis-review}

    \renewcommand{\arraystretch}{1.4}

    \begin{tabularx}{\linewidth}{
        |>{\centering\arraybackslash}p{2.2cm}
        |X
        |>{\centering\arraybackslash}p{1.2cm}|
    }
        \hline
        対象 & \multicolumn{1}{c|}{項目} & 達成度 \\
        \hline

        \multirow{3}{*}{全体}
        & 要求される機能、関係するアクタを網羅している
        & $\bigcirc$ \\
        \cline{2-3}

        & 抽象的な機能要求の記述ではなく、具体的な振る舞いの記述となっている
        & $\bigcirc$ \\
        \cline{2-3}

        & 配置や色、フォントなどが読みやすい
        & $\bigcirc$ \\
        \hline

        \multirow{3}{*}{\shortstack{ユース\\ケース}}
        & 各ユースケースはまとまった振る舞いの単位である
        & $\bigcirc$ \\
        \cline{2-3}

        & 現在時制かつ能動態の動詞句としてユースケース名がつけられている
        & $\bigcirc$ \\
        \cline{2-3}

        & 全てのユースケースの抽象度はおおむね同程度である
        & $\bigcirc$ \\
        \hline

        アクタ
        & 関わるユースケースとのみ関係づけられている
        & $\bigcirc$ \\
        \hline

        \multirow{3}{*}{\shortstack{ユース\\ケース\\記述}}
        & 必要な事前条件、事後条件、基本系列、代替・例外系列を明確かつ簡潔に記述している
        & $\bigcirc$ \\
        \cline{2-3}

        & システムへの入力に加えて、システムからの応答を明記している
        & $\bigcirc$ \\
        \cline{2-3}

        & 各ステップは現在時制と能動句で記述されている
        & $\bigcirc$ \\
        \hline
    \end{tabularx}
\end{table}

\subsubsection*{2.3 システム分析}

%% TODO: BCE（バウンダリ・コントロール・エンティティ）の識別方針を書く

\umlfigure{08a_collaboration_uc1_reserve.png}{1.0}{システム分析のPUMLコラボレーション図（UC1 部屋を予約する）}{collab-uc1}

\umlfigure{08b_collaboration_uc2_checkin.png}{1.0}{システム分析のPUMLコラボレーション図（UC2 チェックインする）}{collab-uc2}

\umlfigure{08c_collaboration_uc3_checkout.png}{1.0}{システム分析のPUMLコラボレーション図（UC3 チェックアウトする）}{collab-uc3}

\umlfigure{08d_collaboration_uc4_inquiry.png}{0.95}{システム分析のPUMLコラボレーション図（UC4 予約内容を確認する）}{collab-uc4}

\umlfigure{08e_collaboration_uc5_cancel.png}{0.95}{システム分析のPUMLコラボレーション図（UC5 予約をキャンセルする）}{collab-uc5}

\umlfigure{09_system_class_diagram.png}{1.0}{システム分析のPUMLクラス図（BCE分類）}{system-class}

\begin{table}[H]
    \centering
    \caption{システム分析のレビュー項目}
    \label{tab:system-analysis-review}

    \renewcommand{\arraystretch}{1.4}

    \begin{tabularx}{\linewidth}{
        |>{\centering\arraybackslash}p{2.2cm}
        |X
        |>{\centering\arraybackslash}p{1.2cm}|
    }
        \hline
        対象 & \multicolumn{1}{c|}{項目} & 達成度 \\
        \hline

        \multirow{3}{*}{全体}
        & ユースケースが要求する全ての振る舞いを網羅している
        & $\bigcirc$ \\
        \cline{2-3}

        & 必要なデータとその置きどころが明確化されている
        & $\bigcirc$ \\
        \cline{2-3}

        & コラボレーション図における相互作用や可能なオブジェクト数が、クラス図の関係・操作・多重度と対応している
        & $\bigcirc$ \\
        \hline

        \multirow{2}{*}{\shortstack{オブ\\ジェクト}}
        & バウンダリ、コントロール、エンティティの3つ組が、それぞれ当該種別の役割に合致して識別されている
        & $\bigcirc$ \\
        \cline{2-3}

        & コントロールはユースケース当たり5個以内である
        & $\bigcirc$ \\
        \hline

        \multirow{3}{*}{相互作用}
        & ユースケースの基本系列に加えて、代替・例外系列も描かれている
        & $\bigcirc$ \\
        \cline{2-3}

        & 描かれた相互作用とその順序が、対応するユースケース記述に合致している
        & $\bigcirc$ \\
        \cline{2-3}

        & バウンダリ、コントロール、エンティティ間で書いてよい相互作用のみが書かれている
        & $\bigcirc$ \\
        \hline
    \end{tabularx}
\end{table}

\subsubsection*{2.4 アーキテクチャ設計}

%% TODO: 採用したアーキテクチャ（多層アーキテクチャ）と選択理由を書く

\umlfigure{10_package_diagram.png}{0.85}{アーキテクチャ設計のPUMLパッケージ図}{package}

\begin{table}[H]
    \centering
    \caption{アーキテクチャ設計のレビュー項目}
    \label{tab:architecture-design-review}

    \renewcommand{\arraystretch}{1.4}

    \begin{tabularx}{\linewidth}{
        |>{\centering\arraybackslash}p{2.2cm}
        |X
        |>{\centering\arraybackslash}p{1.2cm}|
    }
        \hline
        対象 & \multicolumn{1}{c|}{項目} & 達成度 \\
        \hline

        \multirow{5}{*}{\shortstack{アーキ\\テクチャ}}
        & 適切なパターンや実現手段の選択により、全ての非機能要求を満足している
        & $\bigcirc$ \\
        \cline{2-3}

        & 実装上の制約を満足している
        & $\triangle$ \\
        \cline{2-3}

        & 「良い設計のための基本技術」を的確に用いている
        & $\bigcirc$ \\
        \cline{2-3}

        & パッケージ設計原則を満足している
        & $\triangle$ \\
        \cline{2-3}

        & 配置や色、フォントなどが読みやすい
        & $\bigcirc$ \\
        \hline

        \multirow{2}{*}{クラス}
        & 求められる全機能を満足するように操作および属性の割り当てが確定している
        & $\bigcirc$ \\
        \cline{2-3}

        & 実装可能な形で表現されている（型、名前など）
        & $\bigcirc$ \\
        \hline

        \multirow{3}{*}{\shortstack{クラス間\\関係}}
        & 全てについて方向（誘導可能性）が明示されている
        & $\bigcirc$ \\
        \cline{2-3}

        & 全てについて多重度が明示されている
        & $\bigcirc$ \\
        \cline{2-3}

        & 必要な集約、合成集約（コンポジション）の関係が明示されている
        & $\bigcirc$ \\
        \hline
    \end{tabularx}
\end{table}

%% TODO: △とした2項目の理由を本文で補足する
%% ・実装上の制約: 「または他言語」の許容でPythonを選んだため、提供される部分実装には従っていない
%% ・パッケージ設計原則: 6原則のうちCRP・CCPが部分的（層による責務分離とのトレードオフ）

\subsubsection*{2.5 実装}

%% TODO: 実装の構成（言語・フレームワーク・層構成）と、設計モデルとの対応を書く
%% TODO: 必要なら画面のスクリーンショットをここに貼る（\umlfigure は画像なら何でも使える）

\begin{table}[H]
    \centering
    \caption{実装のレビュー項目（抜粋）}
    \label{tab:implementation-review}

    \renewcommand{\arraystretch}{1.4}

    \begin{tabularx}{\linewidth}{
        |>{\centering\arraybackslash}p{2.2cm}
        |X
        |>{\centering\arraybackslash}p{1.2cm}|
    }
        \hline
        対象 & \multicolumn{1}{c|}{項目} & 達成度 \\
        \hline

        \multirow{4}{*}{全体}
        & 設計モデルに適合する形で、構造、振る舞い、関連、相互作用が実装されている
        & $\bigcirc$ \\
        \cline{2-3}

        & 必要に応じて再利用が行われている
        & $\bigcirc$ \\
        \cline{2-3}

        & 意図が分かりやすく実装されている
        & $\bigcirc$ \\
        \cline{2-3}

        & 全ての機能が実装されている
        & $\bigcirc$ \\
        \hline

        \multirow{6}{*}{要素}
        & 全ての属性・変数・引数が適切に初期化・代入されている
        & $\bigcirc$ \\
        \cline{2-3}

        & 各メソッドの大きさと複雑さは適切に抑えられている
        & $\bigcirc$ \\
        \cline{2-3}

        & 命名は分かりやすく一貫している
        & $\bigcirc$ \\
        \cline{2-3}

        & ループ制御の判定が正しい
        & $\bigcirc$ \\
        \cline{2-3}

        & 条件分岐の判定が正しい
        & $\bigcirc$ \\
        \cline{2-3}

        & 扱う事柄に特化した例外を適切に用いている
        & $\bigcirc$ \\
        \hline
    \end{tabularx}
\end{table}

\subsubsection*{2.6 保守}

%% TODO: 保守課題（UC5「予約をキャンセルする」の追加）と、その影響の考察を書く

\begin{table}[H]
    \centering
    \caption{保守課題（UC5の追加）による各工程の差分}
    \label{tab:maintenance-diff}

    \renewcommand{\arraystretch}{1.4}

    \begin{tabularx}{\linewidth}{
        |>{\centering\arraybackslash}p{2.6cm}
        |X|
    }
        \hline
        \multicolumn{1}{|c|}{工程} & \multicolumn{1}{c|}{差分} \\
        \hline

        要求分析
        & UC5をユースケース図に追加。UC4をincludeして予約確認を再利用。ユースケース記述に代替フロー（キャンセル取りやめ）と例外フロー（キャンセル不可状態）を追加 \\
        \hline

        ドメイン分析
        & 予約の状態に「キャンセル済み」を追加 \\
        \hline

        システム分析
        & キャンセル画面・キャンセルCTRLを追加。コラボレーション図を追加。「予約.キャンセルする()」「客室.空室に戻す()」をクラス図に追加 \\
        \hline

        実装
        & キャンセルのコントロール、エンティティのメソッド、route 1本を追加 \\
        \hline
    \end{tabularx}
\end{table}

\subsubsection*{2.7 チェックリストに基づく確認結果}

%% TODO: 全体としてどの程度満たせたかのまとめを書く
%% （レビュー項目は全47項目中45項目が○、2項目が△）

チェックリストによる確認で判明した課題と、その対応結果を表\ref{tab:issues}に示す。

\begin{table}[H]
    \centering
    \caption{確認により判明した課題と対応}
    \label{tab:issues}

    \renewcommand{\arraystretch}{1.3}

    \begin{tabularx}{\linewidth}{
        |>{\centering\arraybackslash}p{0.8cm}
        |X
        |>{\centering\arraybackslash}p{1.8cm}|
    }
        \hline
        \multicolumn{1}{|c|}{\#} & \multicolumn{1}{c|}{課題} & 対応 \\
        \hline

        1 & ドメインクラス図に属性がなく、用語辞書として粗い & 対応済み \\
        \hline
        2 & ドメインオブジェクト図が具体値を持たず「具体的な状況」になっていない & 対応済み \\
        \hline
        3 & ホテル—客室が単純関連（集約が自然） & 対応済み \\
        \hline
        4 & ユースケース図にフロント係→UC4の関連線がなく、UC4の記述と一致しない & 対応済み \\
        \hline
        5 & システム分析クラス図のエンティティ間関連に誘導可能性の矢印がない & 対応済み \\
        \hline
        6 & システム分析クラス図に合成集約がない（予約—宿泊料金／追加サービス利用） & 対応済み \\
        \hline
        7 & パッケージ設計原則 CRP・CCP・SAP が部分的 & 一部対応 \\
        \hline
        8 & 実装上の制約「提供される部分実装に従うこと」に従っていない & 判断を明記 \\
        \hline
        9 & アプリケーション層（コントロール）のテストがない & 対応済み \\
        \hline
        10 & UC2・UC4の代替フロー「氏名・宿泊日で予約を検索する」が未実装 & 対応済み \\
        \hline
        11 & 請求の再取得に対応する操作がシステム分析クラス図・コラボレーション図にない & 対応済み \\
        \hline
        12 & UC5の基本フロー手順5・6（キャンセルの確認）が実装されていない & 対応済み \\
        \hline
    \end{tabularx}
\end{table}

%% TODO: 7と8について、対応しきらなかった理由を本文で補足する
%% ・7: SAPはRepositoryをProtocolで抽象化して充足。CRP・CCPは層による責務分離とのトレードオフとして受容
%% ・8: 「または他言語」の許容に基づく判断であり、代替措置（静的型検査・CI・分析モデルとの対応表）を用意した

\subsubsection*{2.8 その他}

%% TODO: 特記事項があれば書く

\subsection*{3. チーム開発における自身の貢献の説明}

%% TODO: 自身の貢献を書く

\end{document}
